\section{Regional-to-Local Coupling}

The coupling maps Regional2D free-surface elevation and depth-averaged momentum to the Local3D inlet. For a local inward normal \(\mathbf n=(n_x,n_y)\) and tangent \(\boldsymbol\tau=(\tau_x,\tau_y)\),

\begin{equation}
\label{eq:normal-momentum}
q_n=q_xn_x+q_yn_y
\end{equation}

and

\begin{equation}
\label{eq:tangential-momentum}
q_\tau=q_x\tau_x+q_y\tau_y .
\end{equation}

The reconstructed velocities are

\begin{equation}
\label{eq:normal-velocity}
u_n=\frac{q_n}{h}
\end{equation}

and

\begin{equation}
\label{eq:tangential-velocity}
u_\tau=\frac{q_\tau}{h}.
\end{equation}

Discrete discharge preservation is checked by comparing the Regional2D interface flux with the Local3D boundary flux:

\begin{equation}
\label{eq:discharge-preservation}
Q_{2D}
=
\sum_{j=1}^{N_s} q_{n,j}\,\Delta s_j
\approx
\int_{\Gamma_{\mathrm{in}}}\mathbf u\cdot\mathbf n\,\mathrm dA
=
Q_{3D}.
\end{equation}

The G5/G6 prefix-equivalence evidence demonstrates that the accepted G6 forcing prefix is identical to the accepted G5 prefix over 0--600 s.
